\documentclass[border=10pt]{standalone}
\usepackage[siunitx, RPvoltages]{circuitikz}
\usepackage{tikz}

\begin{document}
\begin{circuitikz}[scale=0.8]
    % Operational Amplifier
    \draw (0,0) node[op amp] (opamp) {};
    
    % Input
    \draw (-4,1) node[left] {Input} to[short] (-3,1)
    to[R, l=$R_1$] (-1,1) -- (opamp.+);
    
    % Feedback network
    \draw (opamp.-) to[short] (-1,-1) to[R, l=$R_2$] (3,-1)
    to[switch, l=$S_1$] (5,-1) -- (5,0) -- (opamp.out);
    
    \draw (1,-1) to[R, l=$R_3$] (1,-3) to[switch, l=$S_2$] (5,-3) -- (5,-1);
    
    \draw (-1,-1) to[short] (-1,-5) to[R, l=$R_4$] (3,-5)
    to[switch, l=$S_3$] (5,-5) -- (5,-3);
    
    % Output
    \draw (opamp.out) -- (7,0) node[right] {Output};
    
    % Microcontroller
    \draw[dashed] (7,-6) rectangle (11,-2) node[midway] {Microcontroller};
    \draw (8,-2) -- (8,-1) -| (5.5,-1);
    \draw (9,-2) -- (9,-3) -| (5.5,-3);
    \draw (10,-2) -- (10,-5) -| (5.5,-5);
    
    % Labels
    \node[text width=5cm, align=center, above] at (3.5,3) 
    {Programmable Gain Amplifier for\\Audio Component Testing};
    
    \node[text width=4cm, align=left, below] at (9,-6.5) 
    {S1, S2, S3: Digital\\control switches};
\end{circuitikz}
\end{document}